% =====================================================================
\section{Platform, Build and Running on the RTX 4070}\label{sec:platform}
The source is C++17 with GLFW (window) or EGL (headless), GLAD (OpenGL 3.3 core loader) and no other library. The Makefile targets are \code{all}, \code{test}, \code{docs} (export, tests, workbook), \code{run} and \code{clean};
\code{HEADLESS=1}, \code{CXX}, \code{CC}, \code{OBJ}, \code{APP} and \code{TESTS} select the build.

\begin{table}[h]
\centering\small
\caption{Build routes. ``Verified'' means done in the development environment (Linux/WSL with a Mesa llvmpipe software GPU).}
\label{tab:build}
\begin{tabular}{@{}p{3.1cm}p{7.6cm}p{4.6cm}@{}}
\toprule
Route & Command & State\\\midrule
Windows one click (MinGW-w64) & \code{run.bat} (builds \code{ThermalSense.exe} if missing, then starts it; arguments pass through) & not run here (no Windows machine)\\
Windows make & \code{mingw32-make -j8} (also builds the tests) & not run here\\
Windows cross-build from Linux & \code{tools/cross\_build\_windows.sh} (zig cc, GLFW linked statically from \code{lib/libglfw3.a}) & cross-link only; the executable has never been run natively\\
Linux with a display & \code{./build.sh} or \code{make -j\$(nproc)}, then \code{./run.sh} (needs \code{libglfw3-dev}, \code{pkg-config}) & GLFW route not exercised here\\
Linux headless (EGL only) & \code{./build.sh -{}-headless} or \code{make HEADLESS=1} & verified (this report's tests)\\
CMake & \code{cmake -B build-cmake \&\& cmake -{}-build build-cmake -j} (\code{-DTS\_HEADLESS\_ONLY=ON} for EGL only) & not re-run\\
Tests (this report) & \code{make -j6 HEADLESS=1 CXX=zcxx CC=zcc OBJ=build/obj-e1 APP=build/ts-e1 TESTS=build/tests-e1 build/tests-e1}\newline\code{./build/tests-e1 build/test\_results.json} & verified: \nTests{} tests, \nFail{} failures\\
Exporter & \code{make docs} $\to$ \code{build/scene\_export.json}, \code{build/test\_results.json} & verified\\\bottomrule
\end{tabular}
\end{table}

\paragraph{Test binary.} The tests link only the CORE sources (no \code{src/app}, no \code{src/rendering}, no GL), so simulation, planner,
combat and walk logic are testable without a GPU. Each area has its own file using the \code{TS\_TEST\_GROUP} macro of
\code{tests/TestKit.h}; the groups are listed in Table~\ref{tab:testareas}.

\subsection{Command-line options}
\begin{table}[h]
\centering\small
\caption{Options (\code{src/main.cpp}, \cli{help}).}
\label{tab:cli}
\begin{tabular}{@{}p{6cm}p{9.8cm}@{}}
\toprule
Option & Meaning\\\midrule
(none) & interactive simulation, F1 lists the controls\\
\cli{capture DIR [-{}-only a,b]} & render the report figures (every shot of \code{shotList()}, or the named ones)\\
\cli{record FILE [SECONDS]} & cinematic video, default 120\,s director sequence (needs ffmpeg)\\
\cli{benchmark FILE} & FPS per preset, per-feature ablation and GPU pass times\\
\cli{trace FILE SECONDS} & UAV / rover behaviour log (4\,Hz CSV)\\
\cli{export-docs FILE} & GPU-free export of every part, matrix and mesh (the exporter binary does the same)\\
\cli{quality low|high|ultra|epic} & preset (default: automatic; EPIC on NVIDIA / discrete GPUs)\\
\cli{size WxH}, \cli{fullscreen} & output size; full screen on the primary monitor\\
\cli{headless}, \cli{gpu N} & EGL context without a window; EGL device index (or environment \code{TS\_EGL\_DEVICE})\\
\cli{fps 30|60}, \cli{encoder NAME}, \cli{ffmpeg PATH} & video frame rate, codec (\code{libx264}, \code{h264\_nvenc}, \code{hevc\_nvenc}) and the ffmpeg executable\\\bottomrule
\end{tabular}
\end{table}

\subsection{Running on the RTX 4070 remote machine}
\begin{enumerate}
  \item \textbf{Windows remote.} Install MinGW-w64, run \code{build.bat} (or rebuild the executable with the cross-build script), then
  \code{run.bat -{}-fullscreen} or \code{run.bat -{}-quality epic -{}-size 2560x1440}. For remote desktop use Parsec or Sunshine + Moonlight, which capture the GPU output. Plain RDP often
  exposes only the Microsoft Basic Render Driver (software), so start the program from a console session or use the headless route.
  \item \textbf{Linux desktop.} \code{sudo apt install build-essential pkg-config libglfw3-dev}; \code{./build.sh \&\& ./run.sh -{}-fullscreen}. On hybrid laptops
  \code{\_\_NV\_PRIME\_RENDER\_OFFLOAD=1 \_\_GLX\_VENDOR\_LIBRARY\_NAME=nvidia ./run.sh}.
  \item \textbf{Linux headless over SSH.} Needs only the NVIDIA driver (EGL) and ffmpeg for video; \code{tools/remote\_setup\_linux.sh} installs everything and builds.
\end{enumerate}
The program logs \code{GL\_RENDERER} at start-up. If it says \code{llvmpipe}, \code{Mesa} or \code{Microsoft Basic Render Driver}, the NVIDIA driver is not in use.

\paragraph{FPS on the RTX 4070 has not been measured.} All captures and tests were produced on a Mesa llvmpipe software renderer, where
a shot takes one to two minutes. To obtain the RTX figures run, on the RTX machine:
\begin{center}\small
\code{ThermalSense -{}-headless -{}-quality epic -{}-benchmark build/benchmark.json}\quad (or without \code{-{}-headless} in a window)
\end{center}
and, for the 4K video, \code{ThermalSense -{}-headless -{}-quality epic -{}-size 3840x2160 -{}-fps 60 -{}-encoder h264\_nvenc -{}-record report/demo\_4k.mp4 60}.
The benchmark writes FPS and ms per frame for each preset, per-feature ablation, draw calls and per-pass GPU times; \code{tools/build\_workbook.py} puts them in the workbook.

\paragraph{Capture shots.} \cli{capture} renders the shots of \code{shotList()} at the sizes set in \code{Capture.cpp}. Each shot sets its own time, camera
and flags; the strike shots are rendered at fixed delays after the impact so that each explosion phase is reproducible.
